\section{Attempt 1: Tower Induction in the rocq-coinduction
  library~\cite{pous-companion,schafer-smolka}}
\label{sec:tower}

\subsection{The Tower}

\textbf{``Start at the \code{top} and \code{b} down to the \code{gfp}''}

\begin{lstlisting}
Inductive C: X -> Prop :=
| Cb: forall x, C x -> C (b x)
| Cinf: forall T, T <= C -> C (inf T).
\end{lstlisting}

The tower is closed under \code{b}. The tower is closed under infima (greatest
lower bounds).

Transfiniteness recovered: \code{gfp b = inf C} $= b^{\lambda}(\top)$.

\gap{The unrolling from slides 56--58: \code{inf \{\} = }$\top$, then
\code{b }$\top$, \code{b}$^2\,\top$, \code{b}$^3\,\top$, \ldots,
\code{inf (b}$^n\,\top$\code{)}, converging on \code{inf C = gfp b}. Slide 58
also carries an open problem worth stating in the paper: \emph{prove this is the
same as Parrow \& Weber's ordinal construction}.}

\subsection{The Tower: The best part}

We get an \underline{inductive} proof principle to do \underline{coinduction}!
And indeed to do more\ldots

\subsection{Up-to and coinductive proofs with the tower}

To prove \code{f} is a sound up-to of the tower, \textbf{prove \code{f} holds of
every element of \code{C} by induction on \code{C}.}

This \emph{tower induction} technique is the \underline{single proof method} for
both up-to proofs and coinduction.

Indeed, to prove two elements are related by the \code{gfp}, you can use tower
induction:
\begin{enumerate}
  \item Turn the relation into a predicate \code{P}
  \item Prove \code{P} holds for all elements of the tower by tower induction
  \item The \code{gfp} is in the tower, so \code{P} holds for the \code{gfp}.
\end{enumerate}

This can be done automatically. The base case is almost always trivial.

If we have no base case, we just have to prove stability under \code{b}\ldots
This looks just like coinduction!

\gap{``I can't believe it's not coinduction!'' -- the goal state from slide 63,
where the tower-induction hypothesis is literally the coinduction hypothesis.
This is the payoff of the whole context section and should be a figure.}

\subsection{Tower stats (my claim)}

\gap{The scorecard from slide 64, all five green: 1.~Compositional;
2.~Non-syntactic guardedness; 3.~Nice-to-have: incremental bisimulation;
4.~Amenable to automation (very few tactics, all with uniform behavior);
5.~Smooth way to prove soundness of \textbf{up-to techniques} --- (also) by
\emph{induction!}}
